\begin{helper}[FiberAndDegreeSameColorLiftedIntersectionFixedPairInjectionTail]
\label{helper:FiberAndDegreeSameColorLiftedIntersectionFixedPairInjectionTail}
Let \(m=\noderef{TwoBiteNaturalReducedVertexCount}(n)\), let
\[
p=\noderef{TwoBiteEdgeProbability}(1/2,n),
\qquad L=\log n,
\]
and assume
\[
n\le m^2,\qquad 1\le L,\qquad
n\le 2L^2m .
\]
Fix a red base graph \(G_R\) on \(\operatorname{Fin}(m)\) and distinct
vertices \(r,s\).  If the same-color base common neighborhood satisfies
\[
\left|\{t\in\operatorname{Fin}(m):G_R.Adj(r,t)\ \text{and}\ G_R.Adj(s,t)\}\right|
\le L,
\]
then the number of injections
\[
\phi:\operatorname{Fin}(n)\hookrightarrow
\operatorname{Fin}(m)\times\operatorname{Fin}(m)
\]
for which
\[
\left|
\phi(\operatorname{Fin}(n))\cap
\bigl((N_{G_R}(r)\cap N_{G_R}(s))\times\operatorname{Fin}(m)\bigr)
\right|
>100L^3
\]
is at most
\[
\exp(-2L^3)\,
\left|\{\phi:\operatorname{Fin}(n)\hookrightarrow
\operatorname{Fin}(m)\times\operatorname{Fin}(m)\}\right|.
\end{helper}

\begin{proof}
Put
\[
K=(N_{G_R}(r)\cap N_{G_R}(s))\times\operatorname{Fin}(m)
\subseteq \operatorname{Fin}(m)\times\operatorname{Fin}(m).
\]
The common-neighborhood hypothesis gives
\[
|K|=|N_{G_R}(r)\cap N_{G_R}(s)|\,m\le Lm.
\tag{1}
\]
Thus, for a uniformly chosen image set of size \(n\) in the \(m^2\) product
cells, the number of selected points lying in \(K\) is a hypergeometric
variable with population \(m^2\), marked-set size \(|K|\), and sample size
\(n\).  Its mean is
\[
\mu=n\frac{|K|}{m^2}
\le n\frac{Lm}{m^2}
=L\frac nm
\le 2L^3,
\tag{2}
\]
where the last inequality is the displayed ratio hypothesis.

The injections themselves are ordered samples without replacement.  By
\(\noderef{UniformInjectionImage}\), after quotienting by the number of
orders of a fixed image, counting injections with more than \(100L^3\) marked
image points is equivalent to counting \(n\)-element image subsets with the
same marked-point condition.  Using the standard equivalence between
\(\operatorname{Fin}(m)\times\operatorname{Fin}(m)\) and
\(\operatorname{Fin}(m^2)\), \(\noderef{FinProductTailCountEquiv}\) identifies
this subset count with the hypergeometric tail for the marked subset \(K\).

Apply \(\noderef{HypergeometricCountingTailBound}\) with
\[
N=m^2,\qquad T=100L^3 .
\]
Since \(L\ge1\), \(T>0\).  From (2),
\[
\left(\frac{e\mu}{T}\right)^T
\le
\left(\frac{2eL^3}{100L^3}\right)^{100L^3}
=\left(\frac e{50}\right)^{100L^3}.
\]
Because \(e<3\), the base is at most \(3/50<e^{-1/50}\), and hence the last
quantity is at most \(\exp(-2L^3)\).  Multiplying the resulting probability
bound by the total number of injections gives the claimed cardinality bound.
\end{proof}
